% !TeX program = xelatex
% Geometric meaning of the Legendre transformation.
% Compile: xelatex legendre_transform_tikz.tex
\documentclass[tikz,border=6pt]{standalone}
\usepackage{amsmath}
\usepackage[UTF8,fontset=fandol]{ctex}
\usepackage{fix-cm}
\usetikzlibrary{arrows.meta}

\begin{document}
\begin{tikzpicture}[>=Stealth,thick,every node/.style={font=\fontsize{16}{19}\selectfont}]
  \draw[blue!70!black,->] (-.45,0) -- (6.1,0) node[right] {$\dot q$};
  \draw[blue!70!black,->] (0,-1.95) -- (0,5.85) node[above] {$L$};
  \draw[blue!70!black,domain=.15:5.55,samples=160,smooth]
    plot (\x,{.18*\x*\x});
  \draw[red,domain=.25:5.65,samples=2]
    plot (\x,{1.08*\x-1.62});

  \coordinate (T) at (3,1.62);
  \fill[blue!70!black] (T) circle (2.2pt);
  \draw[densely dashed] (T) -- (3,0) node[below] {$\dot q_*$};
  \draw[densely dashed] (T) -- (0,1.62);
  \fill[red] (0,-1.62) circle (2.2pt);
  \node[red,left] at (0,-1.62) {$-H$};

  \node[blue!70!black,anchor=west] at (5.1,4.6) {$L(\dot q)$};
  \node[red,anchor=west] at (4.5,2.7) {切线};
  \draw[red,->] (4.45,2.6) -- (4.0,2.7);
  \node[green!40!black,align=left,anchor=west,font=\fontsize{14}{17}\selectfont] at (.55,5.15)
    {斜率 $p=L'(\dot q_*)$\\切点与斜率一一对应};
  \node[black!70,align=left,anchor=west] at (.8,-1.35)
    {纵轴截距 $L-p\dot q=-H$};
\end{tikzpicture}
\end{document}
